\documentclass[11pt,a4paper]{article}
\usepackage[utf8]{inputenc}
\usepackage[margin=0.8in]{geometry}
\usepackage{listings}
\usepackage{xcolor}
\usepackage{titlesec}
\usepackage{hyperref}
\usepackage{fancyhdr}

\definecolor{codebg}{RGB}{245,245,245}
\definecolor{codegreen}{RGB}{0,128,0}
\definecolor{codegray}{RGB}{128,128,128}
\definecolor{codeblue}{RGB}{0,0,180}
\definecolor{codepurple}{RGB}{128,0,128}

\lstset{
    language=C++,
    backgroundcolor=\color{codebg},
    basicstyle=\ttfamily\footnotesize,
    keywordstyle=\color{codeblue}\bfseries,
    commentstyle=\color{codegreen}\itshape,
    stringstyle=\color{codepurple},
    numberstyle=\tiny\color{codegray},
    numbers=left,
    numbersep=5pt,
    frame=single,
    breaklines=true,
    breakatwhitespace=false,
    tabsize=4,
    showstringspaces=false,
    captionpos=b,
    morekeywords={override,nullptr,string,shared_ptr,unique_ptr,make_shared,make_unique,
                  lock_guard,unique_lock,mutex,thread,vector,unordered_map,map,
                  chrono,steady_clock,time_point,atomic},
}

\pagestyle{fancy}
\fancyhf{}
\fancyhead[L]{\textbf{LLD --- AWS Payment}}
\fancyhead[R]{\thepage}

\title{\Huge\textbf{AWS Payment Processing}\\[0.3em]\Large Low-Level Design --- C++ Concurrency}
\author{LLD Interview Preparation}
\date{}

\begin{document}
\maketitle
\tableofcontents
\newpage

% ============================================================
\section{File Structure}
% ============================================================
\begin{verbatim}
aws-payment/
  models.hpp       -- PaymentStatus, PaymentMethod, Wallet (walletMtx),
                      PaymentOrder, PaymentEvent
  strategies.hpp   -- PaymentProcessorStrategy, PaymentValidator (Chain),
                      PaymentObserver
  managers.hpp     -- PaymentService (Singleton)
  main.cpp         -- Demo: multiple payment methods, validation,
                      concurrent payments
\end{verbatim}

\textbf{Compilation:}
\begin{verbatim}
g++ -std=c++17 -pthread -o aws_payment.exe main.cpp
\end{verbatim}

\subsection{Concurrency Summary}
\begin{itemize}
    \item \textbf{Level 1:} \texttt{serviceMtx} (PaymentService) --- events, wallets, observers
    \item \textbf{Level 2:} \texttt{walletMtx} (per-Wallet) --- atomic balance check/modify
    \item Released before PSP call (expensive external I/O)
    \item Snapshot pattern for observer notifications
    \item Validator chain immutable after setup --- no lock needed
\end{itemize}

\subsection{Design Patterns}
\begin{itemize}
    \item \textbf{Singleton:} PaymentService
    \item \textbf{Strategy:} CreditCardProcessor, UPIProcessor, WalletProcessor
    \item \textbf{Chain of Responsibility:} AmountValidator, FraudValidator
    \item \textbf{Observer:} AuditObserver, EmailObserver
\end{itemize}

\subsection{Smart Pointers}
\begin{itemize}
    \item \texttt{shared\_ptr<Wallet>} --- service map + WalletProcessor
    \item \texttt{shared\_ptr<PaymentEvent>} --- service map + caller
    \item \texttt{shared\_ptr<PaymentObserver>} --- service vector + caller
    \item \texttt{shared\_ptr<PaymentValidator>} --- chain links
    \item \texttt{shared\_ptr<PaymentProcessorStrategy>} --- caller creates, passes to service
\end{itemize}

% ============================================================
\newpage
\section{models.hpp}
% ============================================================
\begin{lstlisting}[title={\texttt{models.hpp}}]
#ifndef MODELS_HPP
#define MODELS_HPP

#include <string>
#include <vector>
#include <chrono>
#include <iostream>
#include <algorithm>
#include <mutex>
#include <memory>
using namespace std;

// ============================================================
// ENUMS
// ============================================================

enum class PaymentStatus {
    PENDING,
    PROCESSING,
    SUCCESS,
    FAILED,
    REFUNDED
};

enum class PaymentMethod {
    CREDIT_CARD,
    DEBIT_CARD,
    UPI,
    WALLET,
    NET_BANKING
};

inline string statusToString(PaymentStatus s) {
    switch (s) {
        case PaymentStatus::PENDING:    return "PENDING";
        case PaymentStatus::PROCESSING: return "PROCESSING";
        case PaymentStatus::SUCCESS:    return "SUCCESS";
        case PaymentStatus::FAILED:     return "FAILED";
        case PaymentStatus::REFUNDED:   return "REFUNDED";
        default: return "UNKNOWN";
    }
}

inline string methodToString(PaymentMethod m) {
    switch (m) {
        case PaymentMethod::CREDIT_CARD:  return "CREDIT_CARD";
        case PaymentMethod::DEBIT_CARD:   return "DEBIT_CARD";
        case PaymentMethod::UPI:          return "UPI";
        case PaymentMethod::WALLET:       return "WALLET";
        case PaymentMethod::NET_BANKING:  return "NET_BANKING";
        default: return "UNKNOWN";
    }
}

// ============================================================
// WALLET
// Has its own mutex: multiple threads can debit/credit
// the same wallet concurrently (e.g., seller wallet receiving
// credits from different payment events simultaneously).
// PESSIMISTIC: lock before any balance check/modify.
// ============================================================

class Wallet {
    string ownerId;
    double balance;
    vector<string> history;
    mutable mutex walletMtx;

public:
    Wallet(const string& owner, double initial = 0.0)
        : ownerId(owner), balance(initial) {}

    string getOwnerId() const { return ownerId; }

    double getBalance() {
        lock_guard<mutex> lock(walletMtx);
        return balance;
    }

    bool credit(double amount, const string& txnId) {
        lock_guard<mutex> lock(walletMtx);
        if (amount <= 0) return false;
        balance += amount;
        history.push_back("CREDIT: $" + to_string(amount) + " [" + txnId + "]");
        return true;
    }

    bool debit(double amount, const string& txnId) {
        // PESSIMISTIC: lock, check balance, debit -- all atomic
        lock_guard<mutex> lock(walletMtx);
        if (amount <= 0 || balance < amount) return false;
        balance -= amount;
        history.push_back("DEBIT: $" + to_string(amount) + " [" + txnId + "]");
        return true;
    }

    void display() {
        lock_guard<mutex> lock(walletMtx);
        cout << "  Wallet[" << ownerId << "] Balance: $" << balance << endl;
        int start = max(0, (int)history.size() - 5);
        for (int i = start; i < (int)history.size(); i++) {
            cout << "    " << history[i] << endl;
        }
    }
};

// ============================================================
// PAYMENT ORDER
// Represents a single order within a payment event.
// No mutex needed: owned by a single PaymentEvent, modified
// only under the service-level lock.
// ============================================================

class PaymentOrder {
    string orderId;
    string sellerId;
    double amount;
    PaymentStatus status;

public:
    PaymentOrder(const string& orderId, const string& sellerId, double amount)
        : orderId(orderId), sellerId(sellerId), amount(amount),
          status(PaymentStatus::PENDING) {}

    string getOrderId() const { return orderId; }
    string getSellerId() const { return sellerId; }
    double getAmount() const { return amount; }
    PaymentStatus getStatus() const { return status; }
    void setStatus(PaymentStatus s) { status = s; }

    void display() const {
        cout << "    Order[" << orderId << "] Seller: " << sellerId
             << " $" << amount << " " << statusToString(status) << endl;
    }
};

// ============================================================
// PAYMENT EVENT
// Groups multiple orders into one payment transaction.
// shared_ptr<PaymentOrder>: event owns orders, but orders are
// also passed to observers for notification. Shared ownership.
// ============================================================

class PaymentEvent {
    string eventId;
    string userId;
    vector<shared_ptr<PaymentOrder>> orders;
    PaymentMethod method;
    PaymentStatus status;

public:
    PaymentEvent(const string& id, const string& userId, PaymentMethod method)
        : eventId(id), userId(userId), method(method),
          status(PaymentStatus::PENDING) {}

    void addOrder(shared_ptr<PaymentOrder> order) { orders.push_back(order); }

    string getEventId() const { return eventId; }
    string getUserId() const { return userId; }
    PaymentMethod getMethod() const { return method; }
    PaymentStatus getStatus() const { return status; }
    void setStatus(PaymentStatus s) { status = s; }
    const vector<shared_ptr<PaymentOrder>>& getOrders() const { return orders; }

    double getTotalAmount() const {
        double total = 0;
        for (const auto& o : orders) total += o->getAmount();
        return total;
    }

    void display() const {
        cout << "  Event[" << eventId << "] User: " << userId
             << " Method: " << methodToString(method)
             << " Total: $" << getTotalAmount()
             << " Status: " << statusToString(status) << endl;
        for (const auto& o : orders) o->display();
    }
};

#endif
\end{lstlisting}

% ============================================================
\newpage
\section{strategies.hpp}
% ============================================================
\begin{lstlisting}[title={\texttt{strategies.hpp}}]
#ifndef STRATEGIES_HPP
#define STRATEGIES_HPP

#include "models.hpp"
#include <memory>
#include <iostream>
using namespace std;

// ============================================================
// STRATEGY PATTERN: Payment Processor
// unique_ptr<PaymentProcessorStrategy> when only one owner,
// shared_ptr when passed to processPayment (service + caller).
// ============================================================

class PaymentProcessorStrategy {
public:
    virtual ~PaymentProcessorStrategy() = default;
    virtual bool processPayment(shared_ptr<PaymentEvent> event, double amount) = 0;
    virtual string getName() const = 0;
};

class CreditCardProcessor : public PaymentProcessorStrategy {
    string cardNumber;
public:
    CreditCardProcessor(const string& card) : cardNumber(card) {}

    bool processPayment(shared_ptr<PaymentEvent> event, double amount) override {
        cout << "    [CreditCard] Processing $" << amount
             << " on card ****" << cardNumber.substr(cardNumber.length() - 4) << endl;
        // Simulate external PSP call (90% success)
        bool success = (rand() % 10) < 9;
        cout << "    [CreditCard] " << (success ? "OK" : "DECLINED") << endl;
        return success;
    }
    string getName() const override { return "CreditCard"; }
};

class UPIProcessor : public PaymentProcessorStrategy {
    string upiId;
public:
    UPIProcessor(const string& upi) : upiId(upi) {}

    bool processPayment(shared_ptr<PaymentEvent> event, double amount) override {
        cout << "    [UPI] Processing $" << amount << " via " << upiId << endl;
        bool success = (rand() % 10) < 9;
        cout << "    [UPI] " << (success ? "OK" : "TIMEOUT") << endl;
        return success;
    }
    string getName() const override { return "UPI"; }
};

class WalletProcessor : public PaymentProcessorStrategy {
    shared_ptr<Wallet> wallet;  // shared_ptr: wallet also lives in service's map
public:
    WalletProcessor(shared_ptr<Wallet> w) : wallet(w) {}

    bool processPayment(shared_ptr<PaymentEvent> event, double amount) override {
        cout << "    [Wallet] Processing $" << amount
             << " (Balance: $" << wallet->getBalance() << ")" << endl;
        // Wallet::debit is thread-safe (has its own mutex)
        bool success = wallet->debit(amount, event->getEventId());
        cout << "    [Wallet] " << (success ? "OK" : "INSUFFICIENT BALANCE") << endl;
        return success;
    }
    string getName() const override { return "Wallet"; }
};

// ============================================================
// CHAIN OF RESPONSIBILITY: Payment Validation
// shared_ptr<PaymentValidator>: validators are chained via
// shared_ptr links. Service holds the head of the chain.
// ============================================================

class PaymentValidator {
protected:
    shared_ptr<PaymentValidator> next;
public:
    virtual ~PaymentValidator() = default;
    void setNext(shared_ptr<PaymentValidator> n) { next = n; }

    virtual bool validate(shared_ptr<PaymentEvent> event) {
        if (next) return next->validate(event);
        return true;
    }
};

class AmountValidator : public PaymentValidator {
    double maxAmount;
public:
    AmountValidator(double max) : maxAmount(max) {}

    bool validate(shared_ptr<PaymentEvent> event) override {
        double total = event->getTotalAmount();
        if (total <= 0 || total > maxAmount) {
            cout << "    [Validate] FAIL: amount $" << total
                 << " (max $" << maxAmount << ")" << endl;
            return false;
        }
        cout << "    [Validate] OK: amount" << endl;
        return PaymentValidator::validate(event);
    }
};

class FraudValidator : public PaymentValidator {
public:
    bool validate(shared_ptr<PaymentEvent> event) override {
        bool fraud = (rand() % 100) < 5;  // 5% fraud rate
        if (fraud) {
            cout << "    [Validate] FAIL: fraud detected" << endl;
            return false;
        }
        cout << "    [Validate] OK: fraud check" << endl;
        return PaymentValidator::validate(event);
    }
};

// ============================================================
// OBSERVER PATTERN: Payment Notifications
// shared_ptr<PaymentObserver>: shared between service's vector
// and caller (who may want to remove it later).
// ============================================================

class PaymentObserver {
public:
    virtual ~PaymentObserver() = default;
    virtual void onSuccess(shared_ptr<PaymentEvent> event) = 0;
    virtual void onFailure(shared_ptr<PaymentEvent> event, const string& reason) = 0;
    virtual string getName() const = 0;
};

class AuditObserver : public PaymentObserver {
public:
    void onSuccess(shared_ptr<PaymentEvent> event) override {
        cout << "    [Audit] SUCCESS: " << event->getEventId()
             << " $" << event->getTotalAmount() << endl;
    }
    void onFailure(shared_ptr<PaymentEvent> event, const string& reason) override {
        cout << "    [Audit] FAILED: " << event->getEventId()
             << " reason: " << reason << endl;
    }
    string getName() const override { return "Audit"; }
};

class EmailObserver : public PaymentObserver {
    string email;
public:
    EmailObserver(const string& e) : email(e) {}

    void onSuccess(shared_ptr<PaymentEvent> event) override {
        cout << "    [Email -> " << email << "] Payment "
             << event->getEventId() << " succeeded ($"
             << event->getTotalAmount() << ")" << endl;
    }
    void onFailure(shared_ptr<PaymentEvent> event, const string& reason) override {
        cout << "    [Email -> " << email << "] Payment "
             << event->getEventId() << " failed: " << reason << endl;
    }
    string getName() const override { return "Email(" + email + ")"; }
};

#endif
\end{lstlisting}

% ============================================================
\newpage
\section{managers.hpp}
% ============================================================
\begin{lstlisting}[title={\texttt{managers.hpp}}]
#ifndef MANAGERS_HPP
#define MANAGERS_HPP

#include "strategies.hpp"
#include <map>
#include <mutex>
#include <thread>
using namespace std;

// ============================================================
// LOCKING STRATEGY: PESSIMISTIC (local locks, single server)
// ----------------------------------------------------------
// Level 1: serviceMtx (PaymentService)
//   -> protects paymentEvents map, wallet maps, observers vector
//   -> lock_guard for simple reads/writes
//   -> unique_lock in processPayment (release before PSP call
//      and notification -- both are expensive external I/O)
//
// Level 2: walletMtx (per-Wallet)
//   -> protects individual wallet balance
//   -> lock_guard inside debit/credit (atomic check-and-modify)
//   -> Multiple concurrent payments can credit DIFFERENT seller
//      wallets in parallel (no contention)
//
// WHY PESSIMISTIC?
//   Payments involve money -- we can't afford lost updates or
//   double-charges. Lock first, validate, process, commit.
//   Optimistic locking could cause a charge to succeed at PSP
//   but fail the local commit (version mismatch) -- messy refund.
//
// DEADLOCK PREVENTION:
//   - Always Level 1 before Level 2
//   - Release Level 1 before PSP call and notification
//   - notifyObservers takes snapshot under lock, notifies outside
//
// SMART POINTERS:
//   shared_ptr<Wallet>        -> service map + WalletProcessor both hold ref
//   shared_ptr<PaymentEvent>  -> service map + caller both hold ref
//   shared_ptr<PaymentObserver> -> service vector + caller (to remove later)
//   shared_ptr<PaymentValidator> -> chain of responsibility links
//   shared_ptr<PaymentProcessorStrategy> -> caller creates, passes to service
// ============================================================

class PaymentService {
    static mutex singletonMtx;
    mutable mutex serviceMtx;

    map<string, shared_ptr<PaymentEvent>> events;
    map<string, shared_ptr<Wallet>> userWallets;
    map<string, shared_ptr<Wallet>> sellerWallets;
    vector<shared_ptr<PaymentObserver>> observers;
    shared_ptr<PaymentValidator> validatorChain;

    PaymentService() {
        auto amountVal = make_shared<AmountValidator>(10000.0);
        auto fraudVal = make_shared<FraudValidator>();
        amountVal->setNext(fraudVal);
        validatorChain = amountVal;
    }

    // Notify OUTSIDE serviceMtx (snapshot pattern, avoids deadlock)
    void notifySuccess(shared_ptr<PaymentEvent> event) {
        vector<shared_ptr<PaymentObserver>> snap;
        { lock_guard<mutex> lock(serviceMtx); snap = observers; }
        for (auto& obs : snap) obs->onSuccess(event);
    }

    void notifyFailure(shared_ptr<PaymentEvent> event, const string& reason) {
        vector<shared_ptr<PaymentObserver>> snap;
        { lock_guard<mutex> lock(serviceMtx); snap = observers; }
        for (auto& obs : snap) obs->onFailure(event, reason);
    }

public:
    static PaymentService* getInstance() {
        lock_guard<mutex> lock(singletonMtx);
        static PaymentService instance;
        return &instance;
    }

    // --- Observer management (lock_guard: simple scope) ---
    void addObserver(shared_ptr<PaymentObserver> obs) {
        lock_guard<mutex> lock(serviceMtx);
        observers.push_back(obs);
        cout << "[PaymentService] Observer added: " << obs->getName() << endl;
    }

    // --- Wallet management (lock_guard: simple scope) ---
    void createWallet(const string& owner, double balance, bool isSeller = false) {
        lock_guard<mutex> lock(serviceMtx);
        auto w = make_shared<Wallet>(owner, balance);
        (isSeller ? sellerWallets : userWallets)[owner] = w;
        cout << "[PaymentService] " << (isSeller ? "Seller" : "User")
             << " wallet created: " << owner << " ($" << balance << ")" << endl;
    }

    shared_ptr<Wallet> getWallet(const string& owner, bool isSeller = false) {
        lock_guard<mutex> lock(serviceMtx);
        auto& m = isSeller ? sellerWallets : userWallets;
        auto it = m.find(owner);
        return (it != m.end()) ? it->second : nullptr;
    }

    // --- Core: Process Payment ---
    // unique_lock: we release serviceMtx before the PSP call (external I/O)
    // and before notifications. This lets other threads process payments
    // to different sellers in parallel.
    bool processPayment(shared_ptr<PaymentEvent> event,
                        shared_ptr<PaymentProcessorStrategy> processor) {

        cout << "\n=== Processing " << event->getEventId() << " ===" << endl;
        event->display();

        // Step 1: Store event under lock
        {
            lock_guard<mutex> lock(serviceMtx);
            events[event->getEventId()] = event;
            event->setStatus(PaymentStatus::PROCESSING);
        }

        // Step 2: Validate (validator chain is immutable after setup, no lock needed)
        cout << "  [Validation]" << endl;
        if (!validatorChain->validate(event)) {
            event->setStatus(PaymentStatus::FAILED);
            notifyFailure(event, "Validation failed");
            return false;
        }

        // Step 3: Call external PSP -- NO lock held (expensive I/O)
        // This is WHY we use unique_lock pattern: release before I/O
        // so other threads aren't blocked during network calls.
        cout << "  [PSP Call]" << endl;
        double total = event->getTotalAmount();
        bool pspOk = processor->processPayment(event, total);

        if (!pspOk) {
            event->setStatus(PaymentStatus::FAILED);
            for (auto& o : event->getOrders()) o->setStatus(PaymentStatus::FAILED);
            notifyFailure(event, "PSP declined");
            return false;
        }

        // Step 4: Credit seller wallets
        // Each wallet has its own lock (Level 2). No serviceMtx needed.
        // Concurrent payments to DIFFERENT sellers proceed in parallel.
        cout << "  [Wallet Credits]" << endl;
        for (auto& order : event->getOrders()) {
            auto sellerWallet = getWallet(order->getSellerId(), true);
            if (sellerWallet) {
                sellerWallet->credit(order->getAmount(), order->getOrderId());
                order->setStatus(PaymentStatus::SUCCESS);
                cout << "    Credited $" << order->getAmount()
                     << " to " << order->getSellerId() << endl;
            } else {
                order->setStatus(PaymentStatus::FAILED);
                cout << "    Seller wallet not found: " << order->getSellerId() << endl;
            }
        }

        event->setStatus(PaymentStatus::SUCCESS);

        // Step 5: Notify outside lock (snapshot pattern)
        cout << "  [Notifications]" << endl;
        notifySuccess(event);

        cout << "=== " << event->getEventId() << " Complete ===" << endl;
        return true;
    }

    // --- Query operations ---
    void viewEvent(const string& id) {
        lock_guard<mutex> lock(serviceMtx);
        auto it = events.find(id);
        if (it != events.end()) it->second->display();
        else cout << "Event not found: " << id << endl;
    }

    void viewWallet(const string& owner, bool isSeller = false) {
        auto w = getWallet(owner, isSeller);
        if (w) w->display();
        else cout << "Wallet not found: " << owner << endl;
    }

    void listAllEvents() {
        lock_guard<mutex> lock(serviceMtx);
        cout << "\n=== All Payment Events ===" << endl;
        for (auto& [id, ev] : events) {
            ev->display();
            cout << "  ---" << endl;
        }
    }

    PaymentService(const PaymentService&) = delete;
    PaymentService& operator=(const PaymentService&) = delete;
};

mutex PaymentService::singletonMtx;

#endif
\end{lstlisting}

% ============================================================
\newpage
\section{main.cpp}
% ============================================================
\begin{lstlisting}[title={\texttt{main.cpp}}]
#include <iostream>
#include <thread>
#include <vector>
#include "managers.hpp"
using namespace std;

int main() {
    srand(42);  // fixed seed for reproducible demo
    auto* svc = PaymentService::getInstance();

    cout << "=== Payment Processing System ===" << endl << endl;

    // ========== 1. OBSERVER SETUP ==========
    // shared_ptr<Observer>: shared between caller (to remove later)
    // and service's observer vector. Both owners -> shared_ptr.
    auto auditObs = make_shared<AuditObserver>();
    auto emailObs = make_shared<EmailObserver>("user@example.com");
    svc->addObserver(auditObs);
    svc->addObserver(emailObs);

    // ========== 2. WALLET SETUP ==========
    // shared_ptr<Wallet>: stored in service map AND passed to
    // WalletProcessor. Both need it alive -> shared_ptr.
    // Each Wallet has its own mutex for thread-safe debit/credit.
    cout << "\n--- Wallets ---" << endl;
    svc->createWallet("UserA", 500.0);
    svc->createWallet("UserB", 300.0);
    svc->createWallet("UserC", 150.0);
    svc->createWallet("SellerA", 0.0, true);
    svc->createWallet("SellerB", 0.0, true);

    // ========== 3. STRATEGY PATTERN: Different payment methods ==========
    cout << "\n--- Payment 1: Credit Card ---" << endl;
    {
        auto event = make_shared<PaymentEvent>("PAY_001", "UserA", PaymentMethod::CREDIT_CARD);
        event->addOrder(make_shared<PaymentOrder>("ORD_001", "SellerA", 100.0));
        event->addOrder(make_shared<PaymentOrder>("ORD_002", "SellerB", 50.0));
        auto proc = make_shared<CreditCardProcessor>("1234567890123456");
        svc->processPayment(event, proc);
    }

    cout << "\n--- Payment 2: UPI ---" << endl;
    {
        auto event = make_shared<PaymentEvent>("PAY_002", "UserB", PaymentMethod::UPI);
        event->addOrder(make_shared<PaymentOrder>("ORD_003", "SellerA", 75.0));
        auto proc = make_shared<UPIProcessor>("user@paytm");
        svc->processPayment(event, proc);
    }

    cout << "\n--- Payment 3: Wallet ---" << endl;
    {
        auto wallet = svc->getWallet("UserA");
        auto event = make_shared<PaymentEvent>("PAY_003", "UserA", PaymentMethod::WALLET);
        event->addOrder(make_shared<PaymentOrder>("ORD_004", "SellerB", 80.0));
        auto proc = make_shared<WalletProcessor>(wallet);
        svc->processPayment(event, proc);
    }

    // ========== 4. CHAIN OF RESPONSIBILITY: Validation failure ==========
    cout << "\n--- Payment 4: Amount too high (should fail validation) ---" << endl;
    {
        auto event = make_shared<PaymentEvent>("PAY_004", "UserC", PaymentMethod::CREDIT_CARD);
        event->addOrder(make_shared<PaymentOrder>("ORD_005", "SellerA", 15000.0));
        auto proc = make_shared<CreditCardProcessor>("9876543210987654");
        svc->processPayment(event, proc);
    }

    // ============================================================
    // 5. CONCURRENT PAYMENTS
    // ----------------------------------------------------------
    // PESSIMISTIC LOCKING IN ACTION:
    // - processPayment acquires serviceMtx briefly to store event
    // - Releases lock before PSP call (expensive external I/O)
    //   so other threads aren't blocked during network calls
    // - Wallet credits use per-wallet locks (Level 2)
    //   so payments to DIFFERENT sellers proceed in parallel
    // - Notifications use snapshot pattern (copy observers under
    //   lock, notify outside) to prevent deadlock
    // ============================================================
    cout << "\n\n--- Concurrent Payments (3 threads) ---" << endl;
    cout << "All three wallet payments hit at once." << endl;
    cout << "Per-wallet locks prevent double-debit." << endl;
    cout << "Seller wallet credits proceed in parallel." << endl << endl;

    vector<thread> threads;

    auto payFn = [&](const string& eventId, const string& userId,
                     const string& sellerId, double amount) {
        auto wallet = svc->getWallet(userId);
        if (!wallet) return;
        auto event = make_shared<PaymentEvent>(eventId, userId, PaymentMethod::WALLET);
        event->addOrder(make_shared<PaymentOrder>(eventId + "_ord", sellerId, amount));
        auto proc = make_shared<WalletProcessor>(wallet);
        svc->processPayment(event, proc);
    };

    threads.emplace_back(payFn, "CONC_001", "UserA", "SellerA", 50.0);
    threads.emplace_back(payFn, "CONC_002", "UserB", "SellerB", 75.0);
    threads.emplace_back(payFn, "CONC_003", "UserC", "SellerA", 100.0);

    for (auto& t : threads) t.join();

    // ========== 6. VIEW RESULTS ==========
    cout << "\n\n--- Final Wallet Balances ---" << endl;
    svc->viewWallet("UserA");
    svc->viewWallet("UserB");
    svc->viewWallet("UserC");
    cout << endl;
    svc->viewWallet("SellerA", true);
    svc->viewWallet("SellerB", true);

    svc->listAllEvents();

    // ============================================================
    // CONCURRENCY SUMMARY:
    // ----------------------------------------------------------
    // Level 1: serviceMtx (PaymentService)
    //   - lock_guard: wallet create, observer add, event store, queries
    //   - Released before PSP call and notifications
    //
    // Level 2: walletMtx (per-Wallet)
    //   - lock_guard: atomic balance check + modify
    //   - Different wallets lock independently (parallel credits)
    //
    // DEADLOCK PREVENTION:
    //   - Always Level 1 before Level 2 (never reverse)
    //   - Release Level 1 before notifications (snapshot pattern)
    //   - Validator chain is immutable after init (no lock needed)
    //
    // WHY PESSIMISTIC?
    //   Money is involved. Can't afford lost updates or double-charge.
    //   Lock first, validate, process, commit. Simple and correct.
    // ============================================================

    cout << "\n=== Done ===" << endl;
    return 0;
}
\end{lstlisting}

\end{document}
